% Theory 13.8.6

\documentclass{article}

\usepackage{geometry}
\usepackage{amsmath}

% Blackboard letters
\usepackage{amsfonts}

\usepackage{graphicx}
\usepackage{breqn} 
\usepackage{verbatim}
% Hyper links
\usepackage{hyperref}

\title{Theory 13.8.6}
\author{Faisal Memon}

\begin{document}

\maketitle

\section{Aggregation by attention}

The aggregation methods discussed so far either weight the contribution of the neighbors
equally or in a way that depends on the graph topology. Conversely, in graph attention
layers, the weights depend on the data at the nodes. A linear transform is applied to
the current node embeddings so that:

\begin{equation*}
    \mathbf{H}_{k}^{'} = \boldsymbol{\beta}_{k}\mathbf{1}^{T} + \boldsymbol{\Omega}_{k}\mathbf{H}.
\end{equation*}

Then the simularity $ s_{mn} $ of each transformed node embedding $ \mathbf{h}_{m}^{'} $ to the
 transformed node embedding
 $ \mathbf{h}_{n}^{'} $ is computed by concatenating the pairs, taking a dot product
  with a column vector $ \boldsymbol{\phi}_{k} $ of learned parameters, and applying an 
  activation function:
  \begin{equation*}
    s_{mn} = a\left(\boldsymbol{\phi}_{k}^{T}
    \begin{bmatrix}
    \mathbf{h}_{m}^{'} \\ 
    \mathbf{h}_{n}^{'}
    \end{bmatrix}
    \right).
  \end{equation*}

  These variables are stored in a matrix $ \mathbf{S} $ of size $ N \times N $, where each element represents the
  similarity of every node to every other.  As in dot-product self-attention, the attention weights 
  contributing to each output embedding are normalized to be positive and to to one using the softmax operation.
  However, only those values corresponding to the current node and its neighbors should contribute.
  The attention weights are applied to the transformed embeddings:

\begin{equation*}
\mathbf{H}_{k+1} = \mathbf{a}\left[\mathbf{H}_{k}^{'} \cdot \mathbf{Softmask}[\mathbf{S}, \mathbf{A} + \mathbf{I}] \right],
\end{equation*}

where $ \mathbf{a}[\large\bullet] $ is a second activation function.

The function $ \mathbf{Softmask}[\bullet,\bullet] $ computes the attention values by applying softmax
operation separately to each column of its argument $ \mathbf{S} $, but only after setting values where the second
argument $ \mathbf{A} + \mathbf{I} $ is zero to $ -\infty $.  This ensures that the attention to non-neighboring nodes is zero.

This is very similar to the self-attention computation in transformers (see figure 13.12), except that 
(i) The keys, queries, and values are all the same, (ii) The measure of similarity is different,
 and (iii)  The attentions are masked so that each node only attends to itself and its neighbors.
 As in transformers, this system can be extended to use multiple heads that are run in parallel and recombined.
\end{document}